\begingroup
\definecolor{barcategory}{HTML}{A66A19}
\definecolor{barleft}{HTML}{85529B}
\definecolor{barright}{HTML}{285E9A}
\definecolor{baroutput}{HTML}{287B59}
\begin{tikzpicture}[
  x=1cm,y=1cm,font=\small,line width=.65pt,
  arr/.style={postaction={decorate},decoration={markings,
    mark=at position .74 with {\arrow{Stealth[length=4.2pt,width=3.4pt]}}}},
  B/.style={arr,draw=barcategory},
  M/.style={arr,draw=barleft},
  N/.style={arr,draw=barright},
  vertex/.style={circle,fill=black,inner sep=1.8pt},
  lab/.style={fill=white,inner sep=1.5pt,font=\scriptsize}
]

% The two panels have the same vertices and all the same old arrows.
\foreach \panelshift in {0,7.8} {
  \begin{scope}[xshift=\panelshift cm]
    \draw[draw=barcategory!70,dashed,rounded corners=5pt,line width=.55pt]
      (1.4,-.65) rectangle (3.4,1.1);
    \node[fill=white,inner xsep=3pt,inner ysep=2pt,text=barcategory]
      at (2.4,1.1) {$V$};

    \node[vertex,label={[label distance=4pt]90:$a$}] (a) at (0,0) {};
    \node[vertex,label={[label distance=4pt]270:$b$}] (b) at (2.4,0) {};
    \node[vertex,label={[label distance=4pt]90:$c$}] (c) at (4.8,0) {};

    % One middle object and its endomorphism loop are shown.
    \draw[B] (b) .. controls (1.65,.95) and (3.15,.95) .. (b);

    \draw[M] (a) -- (b);
    \draw[N] (b) -- (c);
    \node[text=barleft] at (.8,.3) {$e_{ab}$};
    \node[text=barright] at (4,.3) {$e_{bc}$};
  \end{scope}
}

\node at (2.4,1.7) {$Q^-$};
\node at (10.2,1.7) {$Q^+$};
\draw[right hook-{Stealth[length=6pt,width=4pt]},line width=.8pt]
  (5.7,0) -- node[above=5pt] {$i$} (6.9,0);

% Only the completed quiver has the composite output arrow.
\draw[arr,draw=baroutput,dashed,line width=.9pt]
  (7.8,0) .. controls (8.6,-1.5) and (11.8,-1.5) ..
  node[lab,pos=.5,text=baroutput] {$e_{ac}$} (12.6,0);
\path[use as bounding box] (-.7,-1.35) rectangle (13.3,1.9);
\end{tikzpicture}
\endgroup
